\begin{pastexam}[label={px:ncku-111-cce-ds-01}]{成大電通所 111 學年度 第 1 題（節錄第 b--e 小題）}
（是非）

(20 \textbf{pts}) True or False (一題一分，答錯一題倒扣一分)
\begin{enumerate}[label=\alph*.]
\setcounter{enumi}{1}
\item Recursive function
\begin{enumerate}[label=\Alph*)]
\item Any program that be written using assignment, the if-else statement, and the while
statement can also be written using assignment, if-else, and recursion.
\item Merge sort is stable if we use the iterative procedure to implement it. However, it is not
stable in the recursive version.
\item No matter in the recursive version or in the iterative version, the time complexity of
merge sort is O(n log n), where n is the number of input.
\item Merge sort is an in-place algorithm.
\end{enumerate}

\item Tree
\begin{enumerate}[label=\Alph*)]
\item The time complexity to insert a key in to a binary tree is O(log $n$).
\item The time complexity to increase a value in a Max-Heap is O(log $n$).
\item It will waste $n(k$-1) + 1 fields if we use a list representation to represent each node in
a k-ary tree (i.e., a tree of degree k), where each node uses a list with a fixed size $k$ to
record its children.
\item We need an additional node to represent the overall winner in a loser tree and it runs
slower than a winner tree.
\end{enumerate}

\item Graph
\begin{enumerate}[label=\Alph*)]
\item
\begin{tikzpicture}[baseline=(a.base),v/.style={circle,draw,inner sep=1pt,minimum size=4.5mm,font=\small},>={Stealth[length=1.6mm]}]
\node[v] (a) at (0,0) {2}; \node[v] (b) at (1.1,0) {3};
\draw[->] ([yshift=1.6pt]a.east) -- ([yshift=1.6pt]b.west); \draw[->] ([yshift=-1.6pt]b.west) -- ([yshift=-1.6pt]a.east);
\end{tikzpicture} is a strongly connect component of the graph
\begin{tikzpicture}[baseline=(a.base),v/.style={circle,draw,inner sep=1pt,minimum size=4.5mm,font=\small},>={Stealth[length=1.6mm]}]
\node[v] (a) at (0,0) {1}; \node[v] (b) at (1.1,0) {2}; \node[v] (c) at (2.2,0) {3}; \node[v] (d) at (3.2,0) {4};
\draw[->] ([yshift=1.6pt]a.east) -- ([yshift=1.6pt]b.west); \draw[->] ([yshift=-1.6pt]b.west) -- ([yshift=-1.6pt]a.east);
\draw[->] ([yshift=1.6pt]b.east) -- ([yshift=1.6pt]c.west); \draw[->] ([yshift=-1.6pt]c.west) -- ([yshift=-1.6pt]b.east);
\draw[->] (c) -- (d);
\end{tikzpicture}
\item The precedence relation of the edges in an AOV network can be irreflexive.
\item We can avoid traversing the same edge if the adjacency multilist is used to record a
graph.
\item There exists no backward edge in a breath-first search tree obtained from a undirected
graph.
\end{enumerate}

\item Stack and queue
\begin{enumerate}[label=\Alph*)]
\item The infix expression of the postfix expression ``x 4 + y 7 - / z + '' is
``(x+4)/(y-7)+z''
\item With postfix nation, there is no need for parentheses because there is no ambiguity.
\item Evaluating the postfix expression ``7 2 + 5 8 4 / - *'' returns 27.
\item ``+*abcd'' is a valid prefix expression.
\end{enumerate}
\end{enumerate}
\end{pastexam}

\begin{pastsoln}
原題 a.\ 小題（Objective oriented language 四小句）考物件導向，不在本書範圍，沒有收。 各小題分屬
資結書第~14~章（b）、第~7、8、15~章（c）、第~21--23~章（d）、第~6~章（e），本題放在第~7~章。

答案：
\begin{center}
\begin{tabular}{ccccc}
\toprule
 & A) & B) & C) & D) \\
\midrule
b. & T & F & T & F \\
c. & F & T & T & F \\
d. & F & T & T & T \\
e. & T & T & T & F \\
\bottomrule
\end{tabular}
\end{center}

\textbf{b. Recursive function}
\begin{itemize}
\item A) \textbf{T}。 while 迴圈 \textbf{while} $C$ \textbf{do} $S$ 可以改寫成遞迴函式
$W()$：「\textbf{if} $C$ \textbf{then} $S$；$W()$」，迴圈裡用到的變數當參數傳下去（或放在
共用的變數裡）。 所以 assignment、if-else、while 能寫的程式，assignment、if-else、遞迴也能寫
（反方向就是資結書第~6.4~節的「把遞迴改成迴圈」）。
\item B) \textbf{F}。 合併排序是否穩定只看合併這一步：兩邊的鍵值相同時先取左半邊的，就是穩定的
（資結書第~14.1~、14.4~節）。 遞迴版和迴圈版（由下而上）用的是同一個合併，都是穩定的。
\item C) \textbf{T}。 遞迴版 $T(n) = 2T(n/2) + \Theta(n)$；迴圈版做 $\lceil \lg n \rceil$ 趟、每趟
$\Theta(n)$。 都是 $\Theta(n \log n)$，好、壞、平均都一樣。
\item D) \textbf{F}。 陣列版合併需要 $\Theta(n)$ 的暫存陣列，不是原地排序（資結書第~14.1~節）。
\end{itemize}

\textbf{c. Tree}
\begin{itemize}
\item A) \textbf{F}。 一般的二元樹（或二元搜尋樹）高度可能是 $n - 1$（歪斜），插入要從根走到
葉，最壞 $O(n)$。 只有保證平衡的樹（AVL、紅黑樹）才是 $O(\log n)$。
\item B) \textbf{T}。 最大堆積某個值變大，只可能比父節點大，往上調整（sift-up）最多走樹高
$\lfloor \lg n \rfloor$ 層，$O(\log n)$（資結書第~8.3~節的 increase-key）。
\item C) \textbf{T}。 $n$ 個節點，每個節點固定 $k$ 個子節點欄位，共 $nk$ 個欄位。 除了根，每個
節點恰被一個欄位指到，用到的欄位是 $n - 1$，浪費 $nk - (n - 1) = n(k - 1) + 1$ 個。
\item D) \textbf{F}。 前半句對：敗者樹的內部節點存的是比賽的敗者，總冠軍要另外放在一個節點
（第 $0$ 號）。 後半句錯：輸出冠軍後重新比賽時，敗者樹只要沿著路徑和父節點存的敗者比，不必
去讀兄弟節點，比勝者樹省事，跑得\textbf{比較快}，這正是外部排序用敗者樹的原因（資結書
第~15.3.4~節）。
\end{itemize}

\textbf{d. Graph}
\begin{itemize}
\item A) \textbf{F}。 強連通元件要是\textbf{極大}的互相可達集合（資結書第~23.1~節）。 圖中
$1 \rightleftarrows 2 \rightleftarrows 3$，$1, 2, 3$ 兩兩互相可達；$4$ 沒有出邊，回不去。 所以
強連通元件是 $\{1, 2, 3\}$ 和 $\{4\}$，$\{2, 3\}$ 不是極大的，不是強連通元件。
\item B) \textbf{T}。 AOV 網路的先後關係必須是偏序：遞移而且反自反（irreflexive，沒有活動排在
自己前面），否則就有環、排不出拓樸順序（資結書第~22.4~節）。 所以它可以、而且必須是反自反的。
\item C) \textbf{T}。 鄰接多重串列（adjacency multilist）每條邊只存一個節點，兩端的串列共用它，
節點上有一個標記欄位。 走訪時把看過的邊標記起來，就不會把同一條邊處理兩次（鄰接串列每條邊有
兩份，做不到這件事）。
\item D) \textbf{T}。 無向圖的 BFS 中，非樹邊兩端的層數最多差 $1$，而且一定不是祖先和後代的
關係，所以只有交叉邊，沒有後向邊（資結書第~21.5~節）。
\end{itemize}

\textbf{e. Stack and queue}
\begin{itemize}
\item A) \textbf{T}。 x 4 + 得 $(x+4)$；y 7 $-$ 得 $(y-7)$；$/$ 得 $(x+4)/(y-7)$；z $+$ 得
$((x+4)/(y-7)) + z$。 除法優先於加法，去掉多餘括號就是 $(x+4)/(y-7)+z$。
\item B) \textbf{T}。 後序式的運算順序完全由位置決定，求值時從左到右用一個堆疊就能算，不需要
括號也不需要優先順序（資結書第~6.3~節）。
\item C) \textbf{T}。 7 2 $+$ 得 $9$；8 4 $/$ 得 $2$；5 2 $-$ 得 $3$；$9 \times 3 = 27$。
\item D) \textbf{F}。 合法的前序式運算元要比（二元）運算子多一個。 $+*abcd$ 有 $2$ 個運算子、
$4$ 個運算元，$+(*ab)c$ 已經是完整的式子，最後的 $d$ 多出來，不合法。
\end{itemize}
\end{pastsoln}
